真空中沿 \(z\) 方向有一个尺寸远小于波长 \(\lambda=2\pi c/\omega\) 的振荡电偶极子 \(\boldsymbol p(t)=p_0\cos(\omega t)\hat{\boldsymbol z}\)，可视为理想赫兹偶极子。观测点 \(r\gg\lambda\)，无附近反射体。若把它看作长度为 \(\ell\ll\lambda\) 的短直天线，馈电电流为 \(I(t)=I_0\cos(\omega t)\) 且 \(p_0=I_0\ell/\omega\)。
(1) 写出远区电场和磁场及其极化、相位关系。
(2) 计算单位立体角的平均辐射功率并找出辐射零点。
(3) 积分求总平均辐射功率。
(4) 由输入电流定义辐射电阻，检查低频极限与维度。
